% from ega4-2-fr.tex
\oldpage[IV-2]{33}
schémas quasi-compacts, et pour prouver que $\mathcal L$ est $f$-ample, on
peut appliquer le critère de (II, 4.6.8, c)). Soit donc $\mathcal F$ un
$\mathcal O_X$-Module quasi-cohérent de type fini~; si
\[
\sigma:f^*(f_*(\mathcal F\otimes\mathcal L^{\otimes n}))
\longrightarrow\mathcal F\otimes\mathcal L^{\otimes n}
\]
est l'homomorphisme canonique, on voit comme ci-dessus que $g'^*(\sigma)$ est
l'homomorphisme composé
\[
f'^*(g^*(f_*(\mathcal F\otimes\mathcal L^{\otimes n})))
\xrightarrow{f'^*(u)}f'^*(f'_*(\mathcal F'\otimes
\mathcal L'^{\otimes n}))\xrightarrow{\sigma'}
\mathcal F'\otimes\mathcal L'^{\otimes n}
\]
en posant $\mathcal F'=g'^*(\mathcal F)$, tenant compte de
$(0_{\mathrm I},4.3.3.1)$ et désignant par $u$ l'homomorphisme canonique
\[
g^*(f_*(\mathcal F\otimes\mathcal L^{\otimes n}))
\longrightarrow f'_*(\mathcal F'\otimes\mathcal L'^{\otimes n}).
\]
Or, on sait que $u$ est bijectif quel que soit $n$ (2.3.1)~; d'autre part,
comme $\mathcal F'$ est quasi-cohérent et de type fini, l'hypothèse que
$\mathcal L'$ est ample pour $f'$ entraîne l'existence d'un $n_0$ tel que
pour $n\geq n_0$, $\sigma'$ soit surjectif~; on voit donc que
$g'^*(\sigma)$ est surjectif pour $n\geq n_0$, et comme $g'$ est fidèlement
plat, $\sigma$ est surjectif pour ces valeurs de $n$ (2.2.7), ce qui achève
la démonstration.

\begin{remark}[2.7.3]
\label{IV.2.7.3-fr}
\begin{enumerate}
\item[(i)] Il résulte de (2.6.1), (2.6.4) et (2.5.4.1) que les conclusions
de (2.7.1) sont encore valables dans les cas (i), (iii), (v), (vii) et
(xvi) lorsqu'on suppose seulement que $g$ est quasi-compact et
quasi-fidèlement plat (2.3.3)~; on a déjà remarqué que (2.7.1) est valable
dans le cas (ii) en supposant seulement $g$ surjectif et quasi-compact.

\item[(ii)] Avec les notations et hypothèses de (2.7.1), il peut se faire
que $f$ soit propre et $f'$ projectif sans que $f$ soit quasi-projectif. En
effet, Hironaka [34] a donné un exemple de morphisme propre non projectif
$f:X\to Y$, où $X$ et $Y$ sont deux schémas algébriques réguliers
$(0_{\mathrm I},4.1.4)$ sur un même corps $k$, $Y$ étant projectif sur $k$~;
en outre $Y$ est réunion de deux ouverts affines $Y_i$ $(i=1,2)$ tels que
$f_i:X\times_Y Y_i\to Y_i$ soit projectif pour $i=1,2$. Soit alors
$Y'=Y_1\amalg Y_2$ le préschéma somme~; il est clair que le morphisme
canonique $g:Y'\to Y$ coïncidant avec les injections canoniques dans $Y_1$
et $Y_2$, est fidèlement plat, et il est quasi-compact en vertu de
(I, 5.5.10)~; pourtant, bien que $f':X\times_Y Y'\to Y'$ (coïncidant avec
$f_i$ dans chacun des $Y_i$) soit projectif (II, 5.5.6), il n'en est pas de
même de $f$. Il existe donc un $\mathcal O_{X'}$-Module inversible
$\mathcal L'$ qui est $f'$-ample mais qui n'est pas de la forme
$g'^*(\mathcal L)$ pour un $\mathcal O_X$-Module inversible $\mathcal L$, en
vertu de (2.7.2).

\item[(iii)] Sous les hypothèses de (2.7.1), il peut se faire que $f'$ soit
un isomorphisme local sans que $f$ soit une immersion locale. Soient en
effet $k$ un corps, $\overline k$ une clôture algébrique de $k$, $K$ une
extension séparable de degré fini de $k$, distincte de $k$~; alors le
morphisme structural $f:X\to Y$, où $X=\operatorname{Spec}(K)$ et
$Y=\operatorname{Spec}(k)$, n'est pas une immersion locale, mais si on prend
$Y'=\operatorname{Spec}(\overline k)$, le morphisme $Y'\to Y$ est
fidèlement plat et quasi-compact et $f'$ est un isomorphisme local, puisque
$X'=X\times_Y Y'$ est somme d'un nombre fini de schémas isomorphes à $Y'$.
\end{enumerate}
\end{remark}

\subsection{Préschémas sur une base régulière de dimension 1; adhérence d'un
sous-préschéma fermé de la fibre générique}
\label{subsection:IV.2.8-fr}

\begin{proposition}[2.8.1]
\label{IV.2.8.1-fr}
Soient $Y$ un préschéma localement noethérien, régulier, irréductible et de
dimension 1, de point générique $\eta$, $f:X\to Y$ un morphisme,
$X_\eta=f^{-1}(\eta)$ la fibre au point

